\documentclass{article}
\usepackage{tcee-repaso}
\tcetema[Crecimiento y valoración de empresas]{3.B.4}{Análisis del crecimiento de la empresa. Métodos de valoración de empresas. Especial referencia a los procesos de fusión, adquisición y alianzas estratégicas.}
\tcefecha{08/11/2026}
\begin{document}
\cabecerarepaso
\begin{columnas}

% ==========================================================================
\bloque{Esquema del cante}
\esquemacante{%
\apartadocante{Introducción: «Roma no se hizo en un día»}{2,5}
\apartadocante{I. Teorías del crecimiento de la empresa}{6,5}
\subcante{1. Hasta el tamaño óptimo}{3,5}
\subcante{\hspace*{0.8em}Tecnológico y \autor{Stigler}}{1}
\subcante{\hspace*{0.8em}Transaccional (\autor{Coase}, \autor{Williamson}, GHM)}{2,5}
\subcante{2. Sin tamaño óptimo: \autor{Penrose}, \autor{Gibrat}, \autor{Marris}}{3}
\apartadocante{II. Análisis financiero del crecimiento}{9,5}
\subcante{1. Métodos de valoración}{8}
\subcante{\hspace*{0.8em}Escalera de métodos}{1}
\subcante{\hspace*{0.8em}\autor{Gordon}-\autor{Shapiro}}{1,5}
\subcante{\hspace*{0.8em}Descuento de flujos de caja}{3}
\subcante{\hspace*{0.8em}Múltiplos}{1,5}
\subcante{\hspace*{0.8em}Otros métodos; ¿qué método?}{1}
\subcante{2. Financiación del crecimiento}{1,5}
\apartadocante{III. Estrategias de crecimiento}{10,5}
\subcante{Apertura: matriz de \autor{Ansoff}}{0,5}
\subcante{1. Crecimiento interno}{2}
\subcante{2. Crecimiento externo}{8}
\apartadocante{Conclusión}{1}}
\emph{Pizarra} (solo en I-III): izq., árbol «tamaño óptimo: sí/no» con sus autores; gráfico de gobierno con los modos de crecimiento debajo, (hasta las alianzas). En II, arriba a lo ancho la escalera de métodos; el DCF recuadrado hasta el final; puente $\mathrm{EV}\to V_E$ pequeño. En III, \autor{Ansoff} 2×2 junto al de gobierno; el valor de la adquisición recuadrado al lado del DCF.

% ==========================================================================
\bloque[2,5]{Introducción}
\begin{itemize}
\item \textbf{Enganche}: «Roma no se hizo en un día»: crecer piedra a piedra o comprar la ciudad hecha. OPA de BBVA sobre Sabadell (mayo 2024, hostil, canje, mínimo del 50\,\%): 17 meses discutiendo cuánto valía; aceptación \dato{25,33\,\% del capital}{CNMV, 2025} $\Rightarrow$ sin efecto.
\item \textbf{Relevancia}: el tamaño pesa en productividad, salarios, empleo, resistencia a crisis y competencia. España, país de pymes: \dato{3,31 M de empresas; 54,4\,\% sin asalariados; 2,3\,\% con 20 o más}{INE, DIRCE 2025}; VAB/ocupado de las grandes \dato{1,7 veces el de las micro}{Eurostat, 2024} \vertema{3.A.15}. M\&A mundial \dato{4,6 billones de \$, +49\,\%, 68 operaciones $>$10.000 M}{LSEG, 2025}.
\item \textbf{Contextualización}: micro: la empresa combina factores y maximiza beneficio; ¿por qué existe y dónde acaba? \vertema{3.A.15}. Finanzas: \autor{Friedman} (1970), aumentar beneficios $\to$ \concepto{maximizar el valor para el accionista}. 3 decisiones: inversión \vertema{3.B.2}, financiación \vertema{3.B.3} y \emph{dimensión} (este tema). Ej.: PSA-Opel (2017), Stellantis (2021), CaixaBank-Bankia (2021). \concepto{Crecimiento} = aumento de la dimensión (ventas, activos, empleo, valor); \concepto{interno u orgánico} (invierte en su capacidad), \concepto{externo o inorgánico} (compra capacidad instalada), \concepto{alianza estratégica} (colaboran conservando la personalidad jurídica).
\item \textbf{Problemática}: ¿por qué crecen y hay un tamaño óptimo? ¿Cuánto vale una empresa y por qué difieren los métodos? ¿Crecer dentro, comprar, fusionarse o aliarse? ¿Crean valor las F\&A, y para quién?
\item \textbf{Estructura}: \hyperref[rep:I]{I}, teórico (\emph{por qué}); \hyperref[rep:II]{II}, financiero (\emph{cuánto vale}); \hyperref[rep:III]{III}, estratégico (\emph{cómo}: interno, externo, alianzas = especial referencia). Une organización industrial, finanzas corporativas y economía financiera.
\end{itemize}

% ==========================================================================
\bloque[6,5]{I. Teorías del crecimiento de la empresa\label{rep:I}}
\begin{ideaclave}
No hay una teoría, sino 2 familias: unas ven un tamaño óptimo (tecnología o costes de usar el mercado); otras lo niegan y explican el crecimiento por recursos ociosos, azar u objetivos de los directivos.
\end{ideaclave}

\sub{I.1. Hasta el tamaño óptimo\label{rep:I.1}}
\textbf{a) Enfoque tecnológico (neoclásico)}. La tecnología fija la dimensión eficiente.
\begin{itemize}
\item Rendimientos a escala crecientes (indivisibilidades, especialización) y luego decrecientes (dirección) $\Rightarrow$ CMe a l/p en U \vertema{3.A.12}.
\item \concepto{Escala mínima eficiente} $\overline q$: menor producción con CMe mínimo $\overline c$. Crece hasta $\overline q$: eficiencia \emph{productiva}; la \emph{asignativa} exige $p=\mathrm{CMg}$ (competencia, no tamaño). N.º de empresas $\approx$ demanda a $\overline c$ / $\overline q$ \vertema{3.A.16}.
\item \emph{\autor{Stigler} (1939)}: demanda incierta $\Rightarrow$ planta \concepto{flexible}, CMe algo mayor pero casi constante en un tramo amplio: \emph{horquilla} de tamaños.
\item \emph{Evidencia}: costes en L; \concepto{técnica del superviviente} (\autor{Stigler}, 1958): horquilla amplia.
\item \emph{Límite}: explica la \emph{planta}, no la \emph{empresa}; no dice por qué las economías se aprovechan dentro y no por contrato.
\end{itemize}
\graficorepaso[0.6\linewidth]{costes-l}{Costes medios a l/p en L: $\overline q$ y $\overline c$}

\textbf{b) Enfoque transaccional (\autor{Coase}, \autor{Williamson}, \autor{Grossman}-\autor{Hart})}. \emph{Hacer o comprar} = crecer o no.
\begin{itemize}
\item \autor{Coase} (1937): usar el mecanismo de precios cuesta: \concepto{costes de transacción} \emph{ex ante} (buscar, negociar, redactar) y \emph{ex post} (vigilar, adaptar, resolver). Organizar dentro también cuesta (rendimientos decrecientes de la dirección). Frontera en el \emph{margen}:
  \begin{formula}
  \[\mathrm{CMg}^{\mathrm{org}}(m^{*})=\mathrm{CMg}^{\mathrm{tr}}(m^{*})\]
  \end{formula}
\item \autor{Williamson} (1979, 1985, 1991): racionalidad limitada + oportunismo. 3 atributos: \concepto{especificidad de los activos} $a$ (el clave: cuasirrenta y \emph{hold-up}, monopolio bilateral \vertema{3.A.17}), incertidumbre y frecuencia (no encarece: justifica la estructura propia). Mercado, \emph{híbrida}, jerarquía:
  \begin{formula}
  \[C^{*}(a)=\min\{C_{\text{mercado}}(a),\,C_{\text{híbrida}}(a),\,C_{\text{jerarquía}}(a)\}\]
  \end{formula}
\end{itemize}
\graficorepaso[0.78\linewidth]{gobierno-crecimiento}{Cada estructura de gobierno es un modo de crecer: hilo del tema}
\begin{itemize}
\item $a<a_1$: mercado (subcontratar, no crecer); $a_1<a<a_2$: híbridas (\emph{alianzas}); $a>a_2$: jerarquía (crecer dentro o \emph{comprar}: integración vertical, no indicio de monopolio).
\item \emph{Contratos incompletos} (\autor{Grossman}-\autor{Hart}, 1986; \autor{Hart}-\autor{Moore}, 1990): la propiedad da los \concepto{derechos residuales de control}; integrar da incentivos al comprador y se los quita al vendido $\Rightarrow$ propiedad a quien más aporta; la empresa \emph{no} compra todo.
\item \emph{Valoración}: explica qué se integra, las alianzas y fusiones verticales eficientes; pero da otra frontera de equilibrio.
\end{itemize}

\sub{I.2. Sin tamaño óptimo\label{rep:I.2}}
\begin{itemize}
\item \textbf{\autor{Penrose} (1959)}: la empresa, conjunto de recursos; siempre hay \emph{recursos ociosos} (indivisibilidad, uso múltiple, aprendizaje) que empujan a crecer y se regeneran. \concepto{Efecto \autor{Penrose}}: la dirección limita la \emph{velocidad}, no el tamaño. Diversificación relacionada. Extensión: recursos y capacidades (\autor{Wernerfelt}, 1984). Difícil de medir.
\item \textbf{\autor{Gibrat} (1931)}: \concepto{ley del efecto proporcional}: la tasa de crecimiento es independiente del tamaño y de la historia:
  \begin{formula}
  \[S_{i,t}=(1+\varepsilon_{i,t})S_{i,t-1}\ \Rightarrow\ \ln S_{i,t}\approx\ln S_{i,0}+\textstyle\sum_{\tau=1}^{t}\varepsilon_{i,\tau}\]
  \end{formula}
  $\ln S$ paseo aleatorio $\Rightarrow$ TCL: tamaño \concepto{lognormal} (asimétrica a la dcha.); $\mathrm{Var}(\ln S)\uparrow$ con $t$: dispersión, sin óptimo; concentración por azar. Cola: \concepto{ley de potencias}, \autor{Zipf} $\alpha\approx1$ (\autor{Axtell}, 2001); \autor{Gibrat} + tamaño mínimo la genera (\autor{Gabaix}, 1999). \emph{Evidencia}: vale en grandes; en pequeñas y jóvenes el crecimiento $\downarrow$ con tamaño y edad (\autor{Evans}, \autor{Hall}, 1987; \autor{Sutton}, 1997). Referencia nula; \autor{Geroski} (1999): casi paseo aleatorio (innovaciones irregulares); \autor{Alexander}: crecer para \emph{estabilizar} el beneficio.
\item \textbf{\autor{Marris} (1963-64)}: separación propiedad-control (\autor{Berle}-\autor{Means}, 1932); el directivo maximiza el \emph{crecimiento equilibrado} (demanda = capital) s.\,a. \concepto{restricción de seguridad}: que la \emph{ratio de valoración} (valor de mercado / contable) no caiga tanto como para ser absorbida. Crece más de lo que conviene al accionista. Puentes: ratio $\approx q$ de \autor{Tobin} \vertema{3.A.34}; mercado de control.
\item \emph{Agencia}: \autor{Jensen}-\autor{Meckling} (1976) \vertema{3.A.13}; \autor{Jensen} (1986): \concepto{flujo de caja libre} (lo que sobra tras los VAN $>0$) se gasta en crecer/comprar; la deuda disciplina (LBO). \concepto{Mercado de control empresarial} (\autor{Manne}, 1965): las F\&A, síntoma y remedio.
\end{itemize}

% ==========================================================================
\bloque[9,5]{II. Análisis financiero del crecimiento\label{rep:II}}
\begin{ideaclave}
No hay un valor, sino tantos como preguntas: el contable mira al pasado, el de mercado solo es fiable con eficiencia y el intrínseco descuenta lo que se generará. El precio de una operación es otra cosa: sinergias y negociación.
\end{ideaclave}

\sub{II.1. Métodos de valoración de empresas\label{rep:II.1}}
\begin{itemize}
\item Una \emph{escalera}, de lo que la empresa \emph{tiene} a lo que \emph{generará}: balance $\to$ mercado $\to$ descuento de flujos $\to$ comparables $\to$ mixtos, suma de las partes y opciones reales.
\end{itemize}
\graficorepaso[\linewidth]{metodos-valoracion}{La escalera de métodos (a partir de \autor{Fernández}, 2008)}

\textbf{a) Del valor contable al de mercado}
\begin{itemize}
\item \emph{Balance} (estáticos): \concepto{valor contable} (patrimonio neto = activo $-$ pasivo exigible); \emph{ajustado} (cada partida a mercado); \emph{de liquidación} (suelo); \emph{sustancial} (reposición de los activos de explotación). Dependen de criterios contables y no miran al futuro (marca, clientela, equipo).
\item \concepto{Valor de mercado del capital} = precio $\times$ n.º de acciones. Límites: solo si cotiza; fiable solo con eficiencia semifuerte \vertema{3.B.23}; es un paquete \emph{minoritario}, sin \emph{prima de control}. Precio/valor contable $\approx q$ de \autor{Tobin} = ratio de \autor{Marris}.
\end{itemize}

\textbf{b) Descuento de dividendos: \autor{Gordon}-\autor{Shapiro} (1956)}
\begin{itemize}
\item $P_0=\sum_{t=1}^{N}D_t/(1+r_E)^t+P_N/(1+r_E)^N$, y $P_N$ es el valor de los dividendos posteriores. Supuestos: $D_t=D_0(1+g)^t$, $g<r_E$, retención $b$ reinvertida al $\mathrm{ROE}$. Suma desde $t=1$, razón $(1+g)/(1+r_E)<1$:
  \begin{formula}
  \[P_0=\sum_{t=1}^{\infty}\frac{D_0(1+g)^t}{(1+r_E)^t}=\frac{D_0(1+g)}{r_E-g}=\frac{D_1}{r_E-g},\quad g=b\cdot\mathrm{ROE}\]
  \end{formula}
\item \emph{Implicaciones}: $r_E=D_1/P_0+g$ (coste de los fondos propios) \vertema{3.B.2}; con $D_1=(1-b)\mathrm{BPA}_1$, \emph{PER adelantado}:
  \begin{formula}
  \[\frac{P_0}{\mathrm{BPA}_1}=\frac{1-b}{r_E-b\cdot\mathrm{ROE}};\qquad \frac{\partial P_0}{\partial b}\gtrless0\iff\mathrm{ROE}\gtrless r_E\]
  \end{formula}
  PER $\uparrow$ con crecimiento, $\downarrow$ con riesgo. Retener crea valor solo si $\mathrm{ROE}>r_E$; si $=$, irrelevancia (\autor{Miller}-\autor{Modigliani}, 1961) \vertema{3.B.3}.
\item \emph{Valoración}: sencillo; solo maduras con dividendos estables, muy sensible a $r_E-g$, depende de la política de dividendos $\Rightarrow$ se valora por lo que se \emph{genera}.
\end{itemize}

\textbf{c) Descuento de flujos de caja (DCF)}. El más usado para comprar una empresa entera.
\begin{itemize}
\item \concepto{Flujo de caja libre}: caja de los activos tras impuestos e inversiones, \emph{como si no hubiera deuda} (para accionistas y acreedores) \vertema{3.B.1}:
  \begin{formula}
  \[\mathrm{FCL}_t=\mathrm{EBIT}_t(1-t_{\mathrm{IS}})+\text{amort.}_t-\mathrm{CAPEX}_t-\Delta\mathrm{NOF}_t\]
  \end{formula}
\item \concepto{Coste medio ponderado del capital}, pesos a \emph{mercado}, escudo fiscal en la tasa:
  \begin{formula}
  \[\mathrm{WACC}=r_E(1-L)+r_D(1-t_{\mathrm{IS}})L,\quad L=\frac{V_D}{V},\ V=V_E+V_D\]
  \end{formula}
  $r_E$ por CAPM: $r_f+\beta_E[\mathbb{E}(r_M)-r_f]$ \vertema{3.B.23}; no cotizada: beta de comparable desapalancada y reapalancada. $r_D$: TIR de la deuda neta de gastos.
\item $N$ años con visibilidad + \concepto{valor terminal} (perpetuidad creciente):
  \funcionclave{\mathrm{EV}=\sum_{t=1}^{N}\frac{\mathrm{FCL}_t}{(1+\mathrm{WACC})^t}+\frac{\mathrm{VT}_N}{(1+\mathrm{WACC})^N},\ \ \mathrm{VT}_N=\frac{\mathrm{FCL}_{N+1}}{\mathrm{WACC}-g}}
\item \emph{Puente}: $V_E=\mathrm{EV}-$ deuda neta (deuda financiera $-$ caja y activos líquidos). Se negocia sobre el EV; se paga $V_E$.
\end{itemize}
\graficorepaso[0.92\linewidth]{valor-empresa}{Del valor de la empresa al de las acciones: múltiplos indirectos y directos}
\begin{itemize}
\item \concepto{FCLA} (para el accionista), descontado a $r_E$, da $V_E$ directamente; bien hechos, coinciden:
  \begin{formula}
  \[\mathrm{FCLA}_t=\mathrm{FCL}_t-\text{int.}_t(1-t_{\mathrm{IS}})+\Delta\,\text{deuda neta}_t\]
  \end{formula}
\item \concepto{APV} (\autor{Myers}, 1974): FCL a $r_A$ (sin deuda) + VA del ahorro fiscal $-$ coste de insolvencia; útil si la deuda cambia (LBO). Incertidumbre: escenarios, sensibilidad a WACC y $g$, Monte Carlo.
\item \emph{Valoración}: futuro y caja, no depende de dividendos; pero el VT suele ser la mayor parte y $\mathrm{WACC}-g$ lo mueve mucho: $g\leq$ crecimiento nominal de la economía; contrastar con múltiplos.
\end{itemize}

\textbf{d) Múltiplos (valoración relativa)}
\begin{itemize}
\item \concepto{Múltiplo} = valor de mercado / medida de lo que genera. 2 pasos: calcularlo en comparables (sector, tamaño, crecimiento, riesgo; cotizadas u operaciones) y compararlo con el «teórico».
\item \emph{Indirectos} (EV, luego $-$ deuda neta): EV/EBITDA, EV/EBIT, EV/ventas; no dependen del endeudamiento (EV/EBITDA, tampoco de la amortización). \emph{Directos} ($V_E$): \concepto{PER} $=P/\mathrm{BPA}$, precio/valor contable, rentabilidad por dividendo.
\item Un múltiplo es un DCF resumido: PER alto = crecimiento esperado y poco riesgo, no necesariamente caro. PER: el más usado, recoge la estructura financiera; pero depende del apalancamiento, del ciclo y de la contabilidad, y no sirve con pérdidas. Regla: apalancamientos muy distintos $\Rightarrow$ múltiplos del EV.
\item \emph{Valoración}: rápidos y con precios reales; suponen que el mercado valora bien las comparables. CAPE (\autor{Campbell}-\autor{Shiller}, 1988) para el mercado entero \vertema{3.B.23}.
\end{itemize}

\textbf{e) Otros métodos y ¿qué método?}
\begin{itemize}
\item \emph{Mixtos}: activos + \concepto{fondo de comercio} (lo que el balance no recoge). UEC: capitaliza el \emph{superbeneficio} $B-iV$ durante $n$ años al tipo $k$:
  \begin{formula}
  \[V=A+a_{n,k}(B-iV)\Rightarrow V=\frac{A+a_{n,k}B}{1+i\,a_{n,k}},\ a_{n,k}=\frac{1-(1+k)^{-n}}{k}\]
  \end{formula}
\item \emph{Suma de las partes}: cada negocio con su método, $-$ deuda neta y costes centrales; conglomerados, holdings, inmobiliarias; mide el \emph{descuento de conglomerado}.
\item \concepto{Opciones reales} (\autor{Myers}, 1977): derecho, no obligación, de esperar, ampliar, abandonar o hacer por fases; valen más con más incertidumbre. \emph{VAN ampliado} = VAN estático + opciones (\autor{Trigeorgis}, 1996); esperar tiene valor (\autor{Dixit}-\autor{Pindyck}, 1994) \vertema{3.A.34}. Acciones = \emph{call} sobre los activos con $K$ = deuda (\autor{Merton}, 1974). Se valoran con \autor{Black}-\autor{Scholes} o binomial \vertema{3.B.25}.
\item \emph{¿Cuál?} Ninguno óptimo: DCF para comprar el control o empresas singulares; múltiplos para contraste rápido; balance/suma de partes para inmobiliarias, holdings, liquidación; opciones reales para biotecnología, recursos naturales, \emph{start-ups}. \emph{Start-ups}: múltiplos de ventas o usuarios y \emph{método del capital riesgo}, $V_0=V_N/(1+r_{\mathrm{CR}})^N$, $r_{\mathrm{CR}}$ muy alta.
\item Valor $\neq$ precio: «el precio es lo que pagas; el valor, lo que recibes» (\autor{Graham}, citado por \autor{Buffett}).
\end{itemize}

\sub{II.2. Financiación del crecimiento\label{rep:II.2}}
\begin{itemize}
\item Fondos: beneficios retenidos, socios o deuda. Mercados perfectos: inversión y financiación separadas (\autor{Modigliani}-\autor{Miller}, 1958) \vertema{3.B.3}.
\item Con asimetría y \emph{riesgo moral}, sin fondos propios bastantes no hay crédito aunque el VAN sea $>0$ (\autor{Holmström}-\autor{Tirole}): \emph{ingreso comprometible} $\pi(e^{\mathrm A})(R-B^{\mathrm p}/\Delta\pi)\geq I_0-V_E$ \vertema{3.A.13}.
\item Implicaciones: acumular liquidez; \emph{jerarquía financiera} (\autor{Myers}-\autor{Majluf}, 1984): retenidos $\to$ deuda $\to$ acciones; tecnológicas, capital riesgo; crisis $\Rightarrow$ inversión y compras $\downarrow$; crédito y oleadas de fusiones van juntos.
\end{itemize}

% ==========================================================================
\bloque[10,5]{III. Estrategias de crecimiento\label{rep:III}}
\begin{ideaclave}
Crecer dentro es lento y seguro; comprar es rápido, pero el valor se lo suele llevar el vendedor; aliarse es la vía intermedia cuando los activos son algo específicos. Importa si las sinergias superan la prima.
\end{ideaclave}
\begin{itemize}
\item Cada estrategia = \emph{dirección} + \emph{vía}. Dirección: \concepto{matriz de \autor{Ansoff}} (1957): \emph{penetración} (productos y mercados actuales), \emph{desarrollo de mercados} (productos actuales, mercados nuevos: internacionalización \vertema{3.B.18}), \emph{desarrollo de productos} (nuevos, mercados actuales), \emph{diversificación} (relacionada o conglomerado); riesgo $\uparrow$ hacia la diversificación. Vía: la que dicta $a$ (gráfico de gobierno).
\end{itemize}

\sub{III.1. Crecimiento interno u orgánico\label{rep:III.1}}
\begin{itemize}
\item Proyectos de inversión propios: añade capacidad nueva a la economía. Crece si hay proyectos con VAN $>0$ (rentabilidad $>$ coste del capital): $\mathrm{ROE}>r_E$ \vertema{3.B.2}.
\item \concepto{Crecimiento sostenible} (\autor{Higgins}, 1977): tasa máxima sin cambiar la política financiera. Supuestos: sin acciones nuevas; $b$, endeudamiento (deuda/FP), rotación y margen constantes. FP crecen a $b\cdot\mathrm{ROE}$ $\Rightarrow$ deuda, activo y ventas igual; ROE de \autor{DuPont} \vertema{3.B.1}:
  \funcionclave{g^{*}=b\cdot\mathrm{ROE}=b\cdot\underbrace{\frac{\text{benef.}}{\text{ventas}}}_{\text{margen}}\cdot\underbrace{\frac{\text{ventas}}{\text{activo}}}_{\text{rotación}}\cdot\underbrace{\frac{\text{activo}}{\text{FP}}}_{\text{apalanc.}}}
  Ej.: $0{,}10\cdot2{,}4\cdot5/3=0{,}40$; $b=0{,}5\Rightarrow g^*=20$\,\% (con FP de final de periodo, $g^*=b\,\mathrm{ROE}/(1-b\,\mathrm{ROE})$).
\item Misma $g$ que \autor{Gordon}-\autor{Shapiro}. Crecer $>g^*$: emitir, endeudarse, repartir menos o mejorar margen/rotación; si no, iliquidez: «no estires más el brazo que la manga». Crecer $<g^*$: caja ociosa (FCL de \autor{Jensen}). Une política comercial y financiera.
\item \emph{Valoración}: menos arriesgado, sin prima de control, más barato; pero \emph{lento} (como Roma).
\end{itemize}

\sub{III.2. Crecimiento externo o inorgánico\label{rep:III.2}}
\begin{itemize}
\item Comprar capacidad existente (empresa, división, activos); transfronterizo = IED (frente a \emph{greenfield}) \vertema{3.B.18}. España: \dato{3.336 operaciones, $\approx$103.000 M€}{TTR, 2025}. Concentración, desconcentración y alianzas.
\end{itemize}

\textbf{a) Fusiones y adquisiciones}. Decisión de inversión; la comprada ya tiene precio y hay que pagar más para controlarla.
\begin{itemize}
\item \concepto{Fusión} (RDL 5/2023): 2 o más sociedades en una, que recibe en bloque los patrimonios; \emph{por creación} o \emph{por absorción} (amplía capital). Requisitos: proyecto común, informe de experto (SA/SCA), juntas, 1 mes de oposición de acreedores; canje mal fijado: compensación, no impugnación.
\item \concepto{Adquisición}: compra de acciones o activos hasta el control; ambas conservan personalidad: \emph{grupo}, a menudo bajo \emph{holding}. Cotizada y control: OPA.
\item \emph{Horizontales}, \emph{verticales} o \emph{de conglomerado}.
\item \concepto{Sinergias} $\Delta V=V_{AB}-(V_A+V_B)$; \concepto{prima de control} $P_B-V_B$. Gana el comprador si
  \funcionclave{\underbrace{V_{AB}-(V_A+V_B)}_{\Delta V}-\underbrace{(P_B-V_B)}_{\text{prima}}-\text{costes}>0}
  Cada término, con DCF o múltiplos (\hyperref[rep:II.1]{II.1}); la prima es la parte de las sinergias que se lleva el vendedor; se paga hoy, las sinergias llegan después: «no vendas la piel del oso antes de cazarlo».
\item \concepto{Ecuación de canje} (acciones de $A$ por cada una de $B$, prima $\theta$); $B$ se queda con $\mathrm{TC}\,n_B/(n_A+\mathrm{TC}\,n_B)$; valor razonable y experto independiente:
  \begin{formula}
  \[\mathrm{TC}=(1+\theta)\,\frac{V_{E,B}/n_B}{V_{E,A}/n_A}\]
  \end{formula}
  Ej.: $A$ 100 M acc.\ a 20 €, $B$ 50 M a 10 €, $\theta=30$\,\%, sinergias 300: $\mathrm{TC}=0{,}65$, emite 32,5 M; $B$ tiene 24,5\,\%. Con sinergias (2.800): $B$ +186,8, $A$ +113,2; sin ellas (2.500): $B$ +113,2, $A$ $-$113,2. En efectivo (650): $B$ +150 siempre, $A$ $\pm150$. Acciones: el riesgo se reparte; pagar con acciones, señal de que están caras (\autor{Myers}, 1984).
\item \emph{Motivos que crean valor} (origen de $\Delta V$): costes (escala, alcance, integración vertical: doble marginalización \vertema{3.A.17}; cerrar capacidad); ingresos (ventas cruzadas, recursos de \autor{Penrose}); \emph{poder de mercado} (para el accionista, no para la economía: control); financieros y fiscales; \emph{sustituir a una dirección ineficiente}: mercado de control (\autor{Manne}, 1965).
\item \emph{Que no crean valor}: diversificar (el accionista lo hace solo y más barato) \vertema{3.B.23}; interés de los directivos (\autor{Marris}, \autor{Jensen}); \concepto{arrogancia} (\emph{hubris}, \autor{Roll}, 1986): la prima es transferencia; \emph{maldición del ganador} \vertema{3.A.14}; momento de mercado (\autor{Shleifer}-\autor{Vishny}, 2003); \concepto{adquisiciones asesinas} \dato{5-7\,\% en farmacia de EE.\,UU.}{Cunningham et al., 2021}.
\item \emph{Proceso} (3 fases): diseño (estrategia, candidatas, asesor); ejecución (confidencialidad, carta de intenciones, \emph{due diligence}, negociación, firma, cierre); integración (donde se ganan o pierden las sinergias).
\item \concepto{Compra apalancada} (LBO): fondo $\to$ NEWCO con FP; contrato de financiación (sénior, \emph{mezzanine}) y de adquisición; fusión con la comprada, que paga la deuda con su caja. Deuda \dato{60-90\,\% del precio}{Kaplan-Strömberg, 2009}; MBO si compran los directivos. Valor: disciplina del FCL, vigilancia, escudo fiscal; riesgo de insolvencia. España: fusión tras deuda para comprar en los 3 años previos $\Rightarrow$ plan y experto (RDL 5/2023; asistencia financiera).
\end{itemize}
\graficorepaso[0.6\linewidth]{lbo}{La compra apalancada: NEWCO, financiadores y vendedores}
\begin{itemize}
\item \concepto{OPA}: a todos los accionistas, igual precio; efectivo, canje o mixta (Ley 6/2023; RD 1066/2007). \emph{Obligatoria} al alcanzar el control (30\,\% de voto, o menos si nombra $>$ mitad del consejo en 24 meses), por la totalidad y a precio equitativo; \emph{voluntaria} antes. Amistosa u hostil. \concepto{Deber de pasividad}: el consejo necesita a la junta, salvo buscar \emph{caballero blanco}; reciprocidad. EE.\,UU.: \emph{píldora venenosa}. BBVA-Sabadell: hostil, voluntaria, canje, mínimo 50\,\%: 25,33\,\%: la prima no bastó.
\item \emph{Control de concentraciones}: CNMC (cuota $\geq30$\,\%, salvo adquirida $<10$ M y cuota $<50$\,\%; o $>240$ M en España y 2 con $>60$ M; Ley 15/2007); Comisión (dimensión comunitaria, p.\,ej. $>5.000$ M mundial y $>250$ M en la UE de 2; Reglamento 139/2004). Autorizar, con compromisos o prohibir \vertema{3.B.39} \vertema{4.A.15}.
\item \emph{Evidencia} (estudios de acontecimientos; \autor{Andrade} et al., 2001): comprada gana (prima); compradora poco, nada o pierde; conjunto \dato{1-2\,\% del valor}{Pautler, 2001}. Pérdidas en pocas operaciones muy grandes (\autor{Moeller} et al., 2005); mejor la compra programática.
\end{itemize}

\textbf{b) Desconcentración}. Crecer al revés (descuento de conglomerado). \emph{Escisión total} (se extingue; acciones a los socios); \emph{escisión parcial} (unidad económica; acciones a los \emph{socios}: \emph{spin-off}); \emph{segregación} (acciones a la \emph{propia sociedad}). \emph{Carve-out}, venta de divisiones; las escisiones suelen crear valor (\autor{Cusatis} et al., 1993).

\textbf{c) Alianzas estratégicas}
\begin{itemize}
\item \concepto{Alianza estratégica}: cooperación entre empresas independientes que conservan su personalidad jurídica; \emph{forma híbrida} de \autor{Williamson} (tramo $a_1<a<a_2$).
\item \emph{Societarias} (\emph{joint ventures}): Nokia Siemens Networks (50\,\%, 2007; Nokia la compra en 2013); A220 (Airbus 50,01\,\% en 2018, 75\,\% desde 2020). \emph{Contractuales}: licencias, franquicias, I+D, consorcios; \concepto{UTE} (sin personalidad, obra concreta; Ley 18/1982) y \concepto{AIE} (con personalidad, actividad auxiliar; Ley 12/1991). Participaciones cruzadas minoritarias. \emph{No} son alianzas: grupo/holding (control), subcontratación (mercado), SGR, PERTE.
\item \emph{Ventajas}: rápidas, baratas, flexibles, reparten riesgo, sin prima. \emph{Inconvenientes}: información a un posible competidor, menos control, \emph{hold-up} sin resolver $\Rightarrow$ acaban en compra o ruptura. «Juntos, pero no revueltos».
\item \emph{Competencia}: acuerdo entre competidores puede ser cártel (directrices de cooperación horizontal); \emph{joint venture} con plenas funciones = concentración \vertema{3.B.39}.
\end{itemize}

% ==========================================================================
\bloque[1]{Conclusión}
\begin{itemize}
\item \textbf{Recapitulación}: crecen hasta donde dejan tecnología y costes de transacción, o sin destino (recursos, azar, directivos); valor = VA de lo que se generará; vía según especificidad y sinergias frente a prima. El crecimiento es un mercado en sí mismo.
\item \textbf{Relevancia}: productividad, salarios y competencia; en España, empresas de 250 o más \dato{0,2\,\% del total}{Cifras PYME, 2026}.
\item \textbf{Extensiones}: \vertema{3.A.15} \vertema{3.B.2} \vertema{3.B.3} \vertema{3.B.23} \vertema{3.A.34} \vertema{3.B.18} \vertema{3.B.39} \vertema{4.A.15}; política económica: ciclo, tipos y crédito.
\item \textbf{Opinión}: la adquisición crea valor si lleva los activos a quien mejor los gestiona; desconfiar de sinergias pagadas por adelantado y vigilar el poder de mercado. Mejor quitar regulaciones por umbral de tamaño (\autor{Garicano} et al., 2016; \autor{Almunia}-\autor{López-Rodríguez}, 2018) que subvencionar fusiones.
\item \textbf{Idea final}: Roma no se hizo en un día ni se compra de un día para otro; crecer compensa si lo que se recibe vale más de lo que se paga (\autor{Graham}).
\end{itemize}

% ==========================================================================
\bloque{Autores y fechas}
\hito{1931}{\autor{Gibrat}: efecto proporcional}
\hito{1932}{\autor{Berle}-\autor{Means}: propiedad y control}
\hito{1937}{\autor{Coase}: costes de transacción}
\hito{1939}{\autor{Stigler}: flexibilidad}
\hito{1956}{\autor{Gordon}-\autor{Shapiro}: descuento de dividendos}
\hito{1957}{\autor{Ansoff}: matriz productos-mercados}
\hito{1958}{\autor{Modigliani}-\autor{Miller}; \autor{Stigler}: superviviente}
\hito{1959}{\autor{Penrose}: \emph{The Theory of the Growth of the Firm}}
\hito{1961}{\autor{Miller}-\autor{Modigliani}: dividendos}
\hito{1964}{\autor{Marris}: teoría gerencial}
\hito{1965}{\autor{Manne}: mercado de control}
\hito{1970}{\autor{Friedman}: responsabilidad social = beneficios}
\hito{1973}{\autor{Black}-\autor{Scholes}}
\hito{1974}{\autor{Myers}: APV; \autor{Merton}: acciones como \emph{call}}
\hito{1976}{\autor{Jensen}-\autor{Meckling}: agencia}
\hito{1977}{\autor{Higgins}: crecimiento sostenible; \autor{Myers}: opciones reales}
\hito{1979}{\autor{Williamson}: atributos de la transacción}
\hito{1984}{\autor{Myers}-\autor{Majluf}; \autor{Wernerfelt}: recursos}
\hito{1986}{\autor{Jensen}: FCL; \autor{Roll}: \emph{hubris}; \autor{Grossman}-\autor{Hart}}
\hito{1990}{\autor{Hart}-\autor{Moore}: derechos de propiedad}
\hito{1991}{\autor{Williamson}: formas híbridas; Nobel \autor{Coase} (\autor{Williamson}, 2009; \autor{Hart}-\autor{Holmström}, 2016)}
\hito{1999}{\autor{Gabaix}: \autor{Zipf}; \autor{Geroski}}
\hito{2023}{RDL 5/2023 (modificaciones estructurales); Ley 6/2023 (LMV)}

% ==========================================================================
\bloque{Datos clave}
\begin{itemize}
\item La OPA de BBVA sobre Sabadell la aceptó el \dato{25,33\,\% del capital}{CNMV, 2025}, con un mínimo del 50\,\%.
\item Empresas: \dato{3,31 M; 54,4\,\% sin asalariados; 2,3\,\% con 20+}{INE, 2025}; \dato{0,2\,\% con 250+}{Cifras PYME, 2026}; VAB/ocupado grandes \dato{1,7× micro}{Eurostat, 2024}.
\item F\&A: mundo \dato{4,6 billones de \$ (+49\,\%)}{LSEG, 2025}; España \dato{3.336 operaciones, $\approx$103.000 M€}{TTR, 2025}; capital privado en España \dato{$\approx$7.000 M€, $>$1.000 operaciones}{SpainCap, 2025}.
\item SPAC que salieron a bolsa en EE.\,UU.: \dato{613}{2021}, \dato{57}{2024} y \dato{$\approx$140}{2025; SPACInsider}.
\end{itemize}

% ==========================================================================
\bloque{Preguntas del tribunal}
\begin{enumerate}
\item \textbf{¿Por qué el valor terminal es una perpetuidad creciente y qué exige?} Es \autor{Gordon} aplicado al FCL; exige $g<\mathrm{WACC}$ y $g$ sostenible ($\leq$ crecimiento nominal); suele ser la mayor parte del EV.
\item \textbf{¿Pesos del WACC a valor contable?} No: a mercado (coste de oportunidad actual); deuda tras impuestos, $r_D(1-t_{\mathrm{IS}})$, porque el FCL no recoge el escudo fiscal.
\item \textbf{¿Crea valor retener beneficios?} Solo si $\mathrm{ROE}>r_E$; si $=$, da igual (MM, 1961).
\item \textbf{¿Por qué pierden los compradores?} La prima transfiere al vendedor las sinergias (y más, con \emph{hubris} o puja); se paga hoy y las sinergias son inciertas.
\item \textbf{¿Efectivo o acciones?} Con acciones el riesgo de las sinergias se reparte; con efectivo lo asume el comprador; pagar con acciones puede señalar que están caras.
\item \textbf{¿Cómo se valora una \emph{start-up}?} Múltiplos de ventas o usuarios, capital riesgo ($V_N$ a $r_{\mathrm{CR}}$ alta) y opciones reales.
\item \textbf{¿Cuándo aliarse y no comprar?} Especificidad intermedia ($a_1<a<a_2$): demasiada para el mercado, poca para integrar.
\item \textbf{¿Barreras institucionales al tamaño?} Regulaciones por umbral (50 trabajadores en Francia; 6 M de facturación en España, Grandes Contribuyentes): empresas justo por debajo.
\end{enumerate}

\begin{trampa}
\begin{itemize}
\item \autor{Gordon}: $P_0=D_1/(r_E-g)$ con $D_1=D_0(1+g)$, no $D_0$; suma desde $t=1$; $g=b\cdot\mathrm{ROE}$.
\item WACC: pesos a mercado, $r_D(1-t_{\mathrm{IS}})$; FCL sin intereses (como si no hubiera deuda). En el WACC, $V$ con deuda bruta; en el puente, deuda \emph{neta}.
\item FCLA $=$ FCL $-$ int.$(1-t_{\mathrm{IS}})+\Delta$deuda neta; con 500, 800 al 9\,\%, 25\,\%, +350 y $-$400: 396.
\item \autor{Ansoff}: desarrollo de \emph{productos} = productos nuevos en mercados actuales; de \emph{mercados} = actuales en mercados nuevos.
\item OPA: 30\,\% de voto; comunicación desde el 3\,\%; aceptación 15-70 días naturales; venta forzosa 90\,\% y 90\,\%; blindajes caen al 70\,\%.
\item Escisión parcial: acciones a los socios; segregación: a la propia sociedad.
\item \autor{Gibrat}: tasa \emph{independiente} del tamaño; lognormal es el \emph{tamaño}, no las tasas. Efecto \autor{Penrose}: limita la velocidad, no el tamaño.
\item Frecuencia (\autor{Williamson}): no encarece la transacción. Contratos incompletos: \autor{Grossman}-\autor{Hart}-\autor{Moore}, no \autor{Williamson}.
\end{itemize}
\end{trampa}

\end{columnas}
\end{document}
